% !TEX root = ../thesis.tex

\chapter{绪论}

\section{研究背景与意义}

近年来，低空经济被纳入国家战略发展体系，自 2021 年起首次写入国家规划，2023 年列入战略性新兴产业，2024 年出现在政府工作报告中。无人机作为低空经济的典型载体，已在救援、农业、物流、娱乐等场景中展现出巨大的市场潜力\cite{Cai2014}。本研究课题来源于深圳市科技重大专项"基于数字孪生的城市低空环境仿真建模与动态监测关键技术"（项目编号：KJZD20230923115210021），面向的核心场景为城市低空建筑密集区域，其中由建筑绕流诱发的湍流尾迹与突变风场会对无人机的稳定飞行和推进效率形成严重挑战\cite{Mohamed2023}。

首先，湍流风导致无人机难以稳定飞行：城市突变风场引发强烈的速度脉动与剪切，使无人机推进系统与机体气动负载短时波动剧烈，导致功耗增加、轨迹偏移，甚至触发控制环路振荡而失稳。其次，建筑物干扰叠加形成非均匀、强三维的局部扰动区，随时间快速演化，使传统稳态规划难以补偿，限制低空作业的安全裕度与任务鲁棒性。本质原因在于飞行器对突变气动载荷的响应能力不足，而其中螺旋桨作为推进系统与风场的直接耦合接口，其几何形状对推力输出及扰动抑制具有决定性作用。

在低空经济快速发展的背景下，小型无人机已广泛应用于巡航、救援与配送等任务，但在复杂风场下的稳定性与推进效率仍构成关键瓶颈。回顾已有方法，中低雷诺数螺旋桨气动研究主要包括理论分析、数值模拟与实验测试三大类。理论分析适合前期估算，难以准确获得巡航工况下的细致流场结构；风洞实验虽可信但对流场品质与测量条件要求高、成本与周期较大；相比之下，数值模拟具备几何与工况的高可控性、便于参数化设计，同时具有较好的可重复性与可扩展性，可批量生成训练与验证数据集，为优化算法提供稳定的"状态—反馈"接口以支撑策略学习。

而在数值模拟中，叶素动量理论（Blade Element Momentum Theory, BEMT）虽计算成本低廉，但精度上难以准确捕捉巡航与尾迹耦合；计算流体力学（Computational Fluid Dynamics, CFD）虽精度较高，却在计算成本上难以支撑大规模优化迭代。同时，优化算法虽被广泛用于外形优化，但超参数往往是静态设置，在跨场景任务中表现出明显的收敛困难和鲁棒性不足。仅依靠传统方法难以满足小型无人机在实际复杂风场下的优化需求。

基于上述问题，本文在兼顾数值可靠性与计算效率的前提下，以非线性升力线自由涡迹法（Lifting-Line Free Vortex Wake, LLFVW）为基础，融合机器学习代理建模与智能优化算法开展外形寻优，并在 CFD 中进行高保真验证，由此形成面向小型无人机前飞工况稳定巡航的"建模—优化—验证"一体化方案。

\smallskip

更进一步地，本文聚焦的\textbf{核心科学问题}可归纳为：
\textbf{在缺乏可信度感知的数据驱动代理模型上，气动外形优化算法会被代理在分布外（Out-of-Distribution, OOD）区域的过度自信预测所误导，
收敛到物理不可实现的"虚假最优"；本文研究如何通过对代理模型的不确定度建模、
以及在优化器与代理模型之间引入数据驱动约束，
规避这一陷阱并形成可信的螺旋桨外形优化闭环。}
该问题超出"参数化—代理—优化"的工程链路本身，
直接关联到机器学习代理在工程优化场景中的可信度（trustworthiness）这一前沿议题。

\section{国内外研究现状}

螺旋桨气动外形优化通常包含六个基本环节：（1）确定基准几何模型；（2）设定优化目标函数与关键设计变量；（3）构建性能评估所需的气动仿真模型；（4）匹配优化算法的超参数；（5）在约束条件下开展寻优；（6）对优化结果进行性能验证。本节按螺旋桨外形参数化方法、气动仿真建模技术、气动参数代理建模、外形优化算法、结果验证与可视化五个方面回顾相关研究。

\subsection{螺旋桨外形参数化方法}

对于小型四旋翼无人机，常采用 2$\sim$3 片叶片的布局，典型螺旋桨直径在数十厘米，桨毂直径数毫米。这类全局尺度参数通常在方案期作为约束固定，优化阶段主要关注桨叶沿径向的扭转分布与翼型几何特征。Cho 和 Lee\cite{Cho1995} 将螺旋桨叶片几何分解为扭转角分布、弦长分布与面板节点坐标三参数，并结合涡格法与三维面板法实现优化。
Njaastad 等\cite{Njaastad2022} 以控制点-三次样条插值（Cubic Spline）生成弦长与扭转的径向分布，方法直观、变量集紧凑。
Mola 等\cite{Mola2019} 基于控制点位移、采用三次 NURBS 对弦长、螺距、弯度等轴向分布建模。
Pérez-Arribas 与 Pérez-Fernández\cite{PerezArribas2018} 在 \emph{Advances in Engineering Software} 上提出基于 B-Spline 的螺旋桨叶片设计方法，
以螺旋桨几何的常用设计参数为输入，
通过 B-Spline 曲面 loft 生成叶片的最终表面，
所得到的 B-Spline 表面可直接用于可视化、气动计算与制造工艺，
为本文采用 B-Spline 进行 22 截面径向参数化提供了直接的方法学依据。
张以良等\cite{Zhang2015chinese} 采用三阶贝塞尔曲线对弦长与螺距进行径向拟合并以遗传算法优化控制点，成功重构 4382 型桨、拟合误差低至 0.3\%。综上，B-Spline 兼具光滑性、局部可控性与计算效率的良好平衡，在保持高精度的同时有效压缩变量规模、降低计算复杂度，本文选用 B-Spline 作为螺旋桨外形参数化方案。

\subsection{螺旋桨气动仿真建模技术}

常用的数值计算方法包括升力线理论（Lifting Line Theory, LLT）、涡格法（Vortex Lattice Method, VLM）、叶素动量理论（BEMT）、基于有限体积法的 CFD 等。Goates 等\cite{Goates2022} 基于升力线理论提出一种新型数值升力线方法，能够有效模拟具有任意后掠角、下反角和扭转的机翼气动特性。Boatto 等\cite{Boatto2023} 系统评估了基于叶素动量理论的气动模型在均匀来流下对风电机载荷的预测准确性。涡格法首次出现于 1937 年 Pistolesi 的报告\cite{Pistolesi1937}，Deng 等\cite{DengCJA2021} 提出基于偶极面元的快速多极子方法（Dipole Panel Fast Multipole Method, DPFMM），将非定常涡格法的计算效率提升约两倍。CFD 方法精度较高但计算开销大，Sodja 等\cite{Sodja2012} 在 ANSYS CFX 上对优化螺旋桨开展流场仿真，揭示了优化设计中 CFD 验证的必要性。Ventura Diaz 与 Yoon\cite{Ventura2018} 基于三维非定常 Navier-Stokes 方程，对多种 UAV 平台开展高保真 CFD 仿真。

近年来 LLFVW 方法在风电领域得到广泛应用并为螺旋桨数值模拟提供了新的思路。Marten 等\cite{Marten2016} 在 QBlade 中实现并优化了 LLFVW 模块，
通过自由对流的涡环显式建模尾流，
得到时间精确的诱导速度场，
其计算成本比 CFD 全流场仿真低数个数量级，
适用于在常规动量平衡模型适用范围之外（如大型多兆瓦风机的偏航、塔影、阵风等非定常工况）的气动设计场景，
是本研究 QBlade SIL 求解器内核的算法基础。
Balduzzi 等\cite{Balduzzi2018} 进一步将基于三维非定常 Navier--Stokes 的 CFD 高保真模拟（在 IBM BG/Q 集群上每个算例使用超过 16\,000 核）与 LLFVW 模型相对比，
针对 Darrieus 立轴风力机叶片的非定常气动特性，
验证了 LLFVW 在显式解析尾迹结构与三维效应方面、
能在显著低于 CFD 计算成本的条件下达到工程可用的物理一致性。
两项工作为本研究将 LLFVW 用于螺旋桨的批量样本生成提供了直接的可行性依据。

\subsection{气动参数代理建模}

为节省计算资源，数值模拟中常采用低保真度方法，但这些方法虽然能保持与高保真 CFD 或实验数据相同的趋势，但往往会导致较大的误差。近年来，随着高性能计算和机器学习技术的进步，基于数据驱动的代理模型逐渐成为提高预测效率和精度的重要手段。代理模型通常依赖于大量高质量样本进行训练，这些数据可能来自实验测量或 CFD 仿真。然而，在实际应用中，获取高质量样本的成本较高，并且代理模型对样本分布非常敏感，容易导致过拟合或缺乏泛化能力。

为了克服传统低保真度模型的不足，近年来研究者们提出了多保真度代理模型（Multi-fidelity Models, MFM）。
Bu 等\cite{Bu2022} 提出了一种基于多级层次 Kriging 模型（Multilevel Hierarchical Kriging, MHK）的直升机旋翼气动噪声降低优化方法，
该方法支持三级及以上保真度的数据融合，
在相同噪声降低目标（$-6.84$ dB）下，
单级 Kriging 需要 175 次评估、双级 HK 需要 102 次，
而 MHK 仅需 54 次评估即可达到等效降噪效果，
总成本相比单级 Kriging 减少约 $49\%$。
Chen 等\cite{Chen2022mfm} 提出了一种基于卷积神经网络（Convolutional Neural Network, CNN）的多保真度数据聚合框架 MDA-CNN，
通过构建局部感知域与卷积操作捕捉高保真数据与低保真数据之间的隐式映射关系，
在多尺度仿真、图像融合与实验—仿真混合建模等场景下取得良好适应性。
Wu 等\cite{Wu2024} 提出基于多保真神经网络（Multi-fidelity Neural Network, MFNN）的螺旋桨优化设计框架，
融合低保真 BEMT 与高保真 CFD 数据并结合自适应样本补充策略与教学—学习型优化算法，
将巡航效率从 BEMT 设计下的 $82.3\%$ 提升至 $87.1\%$，
在优化效果与计算开销方面均优于单一保真度方法。
Tao 与 Sun\cite{Tao2019} 在 Aerospace Science and Technology 上提出基于深度信念网络（Deep Belief Network, DBN）作为低保真模型的线性回归多保真度框架，
针对马赫数不确定度下的翼型与机翼鲁棒优化，
显著提升了优化收敛效率。
Mian 等\cite{Mian2021} 采用空间映射（Space Mapping）方法、
结合物理基础代理模型对小型电动螺旋桨开展设计优化，
通过引入低保真度与高保真度模型之间的映射关系避免对高保真度模型直接寻优，
成功降低了计算成本。

尽管代理模型在多学科优化中取得了显著应用，但在无人机螺旋桨设计中，传统叶素动量法在巡航状态下依赖经验公式修正，导致结果不够精确；CFD 仿真虽然能够提供高精度的气动性能预测，但其计算成本高昂，难以进行大规模优化。为了解决这些问题，LLFVW 方法被引入，以更精确地仿真巡航角度和尾流效应，并通过深度神经网络（DNN）作为代理模型，利用螺旋桨气动模拟数据的近似，显著降低计算成本。

\subsection{螺旋桨气动外形优化算法}

传统的数值优化方法（如最速下降法、拟牛顿法、共轭梯度法等）虽然在局部优化问题中表现出色，但由于其对梯度信息的依赖，往往容易陷入局部最优解，且问题的适应性差，难以解决大规模、多目标或复杂约束的优化问题。近年来人工智能和机器学习技术的快速发展催生了许多新型智能优化算法，如遗传算法、粒子群优化算法、模拟退火算法等。

遗传算法是一种应用比较广且开始时间比较早的优化算法。Shahrokhi 等\cite{Shahrokhi2007}、Jing Zhang 等\cite{Zhang2021}、Jiaqi Liu 等\cite{Liu2023} 分别基于遗传算法对翼型升阻比进行优化。粒子群算法方面，刘坤澎等\cite{LiuKP2023} 以某太阳能无人机为例设计了一种能够满足多设计点螺旋桨气动外形优化和多工况变桨距的方法。

Ziyang Liu 等\cite{Liu2024} 采用深度强化学习（DRL）方法，成功地最大化了翼型的升力系数，从而有效解决了超临界机翼的优化设计问题。西北工业大学邬晓敬等\cite{Wu2022} 提出了一种基于"教"与"学"概念的优化算法，在 200 维高维优化问题上表现出较强的适应性和高效性。

在多目标优化方向，Oliveira 等\cite{Oliveira2023} 采用遗传算法对螺旋桨进行多目标优化，得到 Pareto 前缘。
Park 等\cite{Park2022} 也采用基于遗传算法的多目标优化方法，对火星探测器的螺旋桨进行了形状优化。
针对多旋翼平台，Klimczyk\cite{Klimczyk2022} 在 \emph{Aircraft Engineering and Aerospace Technology} 上提出了一种基于叶素法和演化算法的螺旋桨自定义设计方法，
在叶片厚度、前后缘加工约束下进行参数化优化，
并使用 URANS 仿真在 hex-rotor 构型下评估螺旋桨在悬停工况的载荷分布，
但其优化目标局限于悬停推力与扭矩，
未涉及前飞工况下前后桨差速配平问题。
针对低空经济场景的城市巡航需求，
本文将优化对象由悬停工况扩展至前飞工况，
并将配平方程显式嵌入到优化目标函数评估流程中（详见第~\ref{chap:optimization}~章）。

为提升算法自适应性，研究者开始将强化学习机制引入进化优化过程。Chen 等\cite{Chen2025rl}、Liu 等\cite{Liu2024rl}、Sakurai 等\cite{Sakurai2010rl} 分别在灰狼优化、多目标差分进化和遗传算法中嵌入了 RL 控制。Li 等\cite{Li2024review} 综述指出 "RL 控制进化参数" 是一条具有良好适应性与可推广性的路径。然而，无论是确定性进化算法还是 RL 调参的进化算法，其搜索可靠性均建立在所依赖代理模型的可信度之上；当代理模型在 OOD 区域过度自信时，优化算法的差异并不能纠正由代理引发的误导。该结论也得到了本研究在第~\ref{chap:surrogate}~章 V1 阶段以 PPO + NSGA-II 与 CMA-ES 双路线对照实验的支持。

\subsection{优化结果验证与可视化分析}

当前主流研究普遍采用 CFD 平台内置的可视化模块（如 Fluent-Post、ParaView、Tecplot 等）对仿真数据进行多维呈现，途径涵盖性能参数曲线（如推力系数 $C_T$、功率系数 $C_P$）、速度矢量图、压力等值图、涡量场、$Q$-criterion 涡结构识别等多种方式。这些可视化手段不仅便于定量对比桨叶设计优化效果的差异，更能够以可视化的方式揭示关键气动现象的物理机制。

\section{研究场景与工作范围}
\label{sec:scope}

本文研究场景明确定位为：
\textbf{恒定自由来流（$V_{\text{cruise}}=10$\,m/s）前飞工况下、
满足三轴差速配平约束的 X 构型小型四旋翼无人机螺旋桨气动外形优化}。
选择前飞工况作为研究场景的核心理由如下：

\begin{enumerate}
  \item \textbf{工程价值明确。}
        前飞巡航是小型无人机在物流配送、城市低空巡航等应用场景下耗能占比最大、推进效率影响最显著的飞行阶段；
        而现有螺旋桨气动外形优化研究多集中于悬停工况
        （如 Klimczyk\cite{Klimczyk2022} 等），
        前飞巡航工况下的优化研究相对稀缺。
  \item \textbf{前飞工况本身具有足够的气动复杂性。}
        恒定自由来流下的 X 构型四旋翼前飞已经引入了\emph{非轴向来流的尾迹—桨叶耦合}、
        \emph{前后桨升力差导致的俯仰矩耦合}、
        \emph{三轴差速配平约束}等典型复杂气动现象，
        构成对代理模型可信度与优化器—代理耦合的有意义考核场景，
        为本文方法学链条的验证提供了足够的复杂性。
  \item \textbf{方法学先行。}
        在代理模型可信度（第~\ref{chap:surrogate}~章）与优化器—代理耦合（第~\ref{chap:optimization}~章）
        在恒定自由来流工况下完整闭环验证之后，
        将湍流风场扰动作为后续方法学的扩展（毕业论文阶段）即可获得增量价值；
        反之若不先解决稳态工况下的代理可信度问题，
        直接叠加湍流风场只会放大 OOD 现象、使方法学难以解耦验证。
\end{enumerate}

因此，本中期论文的工作范围严格限定在上述\emph{恒定自由来流前飞工况}下；
湍流风场建模与风场—优化耦合作为毕业论文阶段的扩展工作，
不在本中期论文范围内。

\section{研究内容与技术路线}

在上述工作范围下，
针对开题报告所提出的"建模—优化—验证—可视化"一体化方案，
并结合中期阶段实际研究进展中发现的 OOD 外推陷阱问题，
本文的研究内容确定为以下四项：

\begin{enumerate}
  \item \textbf{螺旋桨气动外形参数化与 LLFVW 自动化仿真平台搭建。}
        采用 8 个 B-Spline 控制点（4 chord + 4 twist, degree=3）参数化生成 22 个径向截面，
        基于 QBlade SIL 接口构建跨平台的自动化仿真框架，
        通过 UIUC 风洞数据库完成翼型与螺旋桨两级验证。
  \item \textbf{基于 MoE 的不确定度量化代理模型构建。}
        在单 MLP DNN 基线（V1）暴露 OOD 虚假最优问题后，
        将代理模型升级为 $K=4$ 混合专家网络，
        采用异方差负对数似然损失训练，
        同步输出预测均值 $\mu$ 与不确定度 $\sigma$。
  \item \textbf{CMA-ES 配合数据驱动硬约束的配平嵌套优化。}
        将原 PPO 在线调参的 NSGA-II 调整为 CMA-ES，
        新增三轴配平嵌套求解器（求解 $\alpha$、$\Omega_1$、$\Omega_2$），
        通过控制点范围硬约束消除 OOD 风险，
        形成 $\sigma$ 引导的优化—主动补点闭环。
  \item \textbf{CFD 高保真验证模板的搭建。}
        搭建 Star-CCM+ CFD 验证模板（$\sim 700$ 万网格、$y^+ \in [0,1]$、阻塞比 2.5\%、SST $k$-$\omega$），
        在 APC 10$\times$7 基线桨上完成模板验证（$C_P$、$C_T$、$\eta$ 误差 $\leq 5\%$），
        作为优化最终解的独立高保真验证手段，
        与 QBlade LLFVW 同源回验形成"低保真高频闭环 + 高保真终验"的双层验证体系。
\end{enumerate}

\section{论文结构安排}

本文剩余章节安排如下：

\textbf{第~\ref{chap:simulation}~章}~~介绍螺旋桨气动外形 B-Spline 参数化建模、LLFVW 气动求解器、自动化仿真框架以及基于 UIUC 数据库的两级验证。

\textbf{第~\ref{chap:surrogate}~章}~~阐述代理模型从单一 MLP（V1）演进至混合专家（MoE, V2）的过程，重点剖析 V1 在优化过程中暴露的 OOD 虚假最优问题及 MoE 异方差不确定度量化的解决方案。

\textbf{第~\ref{chap:optimization}~章}~~阐述 CMA-ES 配合数据驱动硬约束的优化框架与三轴配平嵌套求解器，给出 V1$\to$V4 四轮迭代的对比结果。

\textbf{第~\ref{chap:wind_cfd}~章}~~给出基于 Star-CCM+ 的 CFD 高保真验证模板设置（计算域、网格、$y^+$、湍流模型）以及"QBlade 同源回验 + CFD 终验"的双层验证策略。

\textbf{第~\ref{chap:plan}~章}~~汇报当前 1000 几何全量数据生成进展、与开题计划的对比，以及后续研究安排。
